\documentclass[11pt]{article}

\usepackage[margin=1in]{geometry}
\usepackage{amsmath,amssymb,amsthm}
\usepackage{bm}
\usepackage{hyperref}
\usepackage{graphicx}
\usepackage{booktabs}
\usepackage{enumitem}
\usepackage{courier}
\usepackage{listings}

\title{Temporal Stake Bonding for Governance of Cross-Organizational Calibration Signals}

\author{The Phase Mirror Research Initiative \\ \\ Citizen Gardens \\ The Foundation of Multiplicity}
\date{April 26, 2026}

\begin{document}
\maketitle

\begin{abstract}
We propose a temporal stake bonding mechanism and a three-phase cold-start policy for governance of cross-organizational calibration signals.
Unlike slashing schemes developed for blockchain consensus, our setting is a shared false-positive (FP) calibration network that tunes governance rules across multiple organizations.
The core ideas are:
(i) viewing \emph{time-of-detection} of misbehavior as a first-class variable in stake slashing for governance, not just for consensus, and
(ii) applying a three-phase cold-start regime (probation, graduated weight, full weight) to reputation-based aggregation of calibration signals from cross-organizational nodes.
We also give a qualitative mapping from these mechanisms to a Phase Mirror-style FP calibration oracle, while explicitly leaving many implementation details---stake levels, currency units, subscription models, revenue-sharing fractions---out of scope to preserve design flexibility and trade-secret space.
\end{abstract}

\newpage
\tableofcontents
\newpage
\section{Introduction and Motivation}

Cross-organizational calibration networks arise when multiple organizations contribute labeled or scored events to maintain shared governance rules, such as thresholds for false-positive (FP) handling, anomaly detection, or other policy decisions.
In such networks, the problem is not payment settlement or ledger consensus, but rather the governance of \emph{calibration signals}: which validators should have influence on shared thresholds, and how should misbehavior be penalized.

Na{\"i}ve adaptations of blockchain-style slashing to this setting have two immediate issues:
\begin{itemize}[leftmargin=*]
  \item \textbf{Temporal insensitivity.}
  Traditional slashing often treats misbehavior as a static event: once detected, a fixed fraction of stake is removed, regardless of how long the misbehavior went undetected.
  In a governance network, however, late-detected misbehavior can have much higher impact on downstream decisions, suggesting that \emph{detection delay} should matter.
  \item \textbf{Cold-start symmetry.}
  If all nodes begin with identical reputation and weight, a newly joined or under-observed organization may wield the same influence as a long-standing, well-behaved one during the early stages.
  This creates attack surfaces where adversarial nodes can exploit the cold-start period before their behavior is sufficiently observed.
\end{itemize}

This paper addresses both issues through:
\begin{enumerate}[leftmargin=*]
  \item A \emph{temporal stake bonding} model where penalties for misbehavior grow with the delay between event emission and detection, and where rapid \emph{self-reporting} of errors yields reduced penalties.
  \item A \emph{three-phase cold-start regime} (probation, graduated weight, full weight) for governance nodes in a cross-organizational calibration system, with illustrative parameter choices for submission-count thresholds.
\end{enumerate}

Our aim is to freeze these concepts as prior art for governance and calibration networks---particularly Phase Mirror-style FP calibration oracles---without binding any one implementation to specific stake levels, currencies, pricing models, or detailed revenue-sharing schedules.

\section{System Overview}

\subsection{Cross-Organizational Calibration Network}

We consider a network of organizations (validators) that contribute calibration signals for shared governance rules.
Each validator submits events of the form:
\[
e = (v, r, \ell, t_0),
\]
where:
\begin{itemize}[leftmargin=*]
  \item \(v\) identifies the validator (e.g., an organizational node).
  \item \(r\) identifies a governance rule or calibration target (for example, a threshold for flagging FP cases).
  \item \(\ell\) encodes a label or signal (e.g., ``FP'' vs ``non-FP'').
  \item \(t_0\) is the event time.
\end{itemize}

A separate process---using aggregate statistics, human review, or external ground truth---later determines whether a given event corresponds to correct behavior or misbehavior (e.g., intentional corruption of FP labels or disregard of agreed-upon calibration protocols).
Misbehavior is detected at some time \(t_{\mathrm{detect}}>t_0\), yielding a detection delay \(d = t_{\mathrm{detect}} - t_0\).

\subsection{Stake and Reputation in Governance Networks}

Each validator stakes an amount \(S_v \ge 0\) in some unit (not specified here).
The network maintains a reputation score \(R_v\) that reflects historical correctness and participation quality.
Unlike pure consensus protocols, we treat stake and reputation as tools for \emph{governance of calibration signals}, not for ordering transactions or preventing double-spends.

The influence of a validator on shared rules is determined by a \emph{weight} \(w_v\) that depends on both:
\begin{itemize}[leftmargin=*]
  \item its current reputation \(R_v\), and
  \item its cold-start phase (Section~\ref{sec:cold-start}).
\end{itemize}

\section{Temporal Stake Bonding}
\label{sec:temporal-bonding}

\subsection{Qualitative Model}

In standard slashing, misbehavior triggers a penalty that is often characterized solely by the type or count of infractions and the current stake, without explicit regard for how long it took to detect those infractions.
We instead treat \emph{time-of-detection} (or detection delay) as a first-class variable:

\begin{itemize}[leftmargin=*]
  \item For each misbehavior event, the system computes a detection delay \(d\) between the event time and the time at which the misbehavior is confirmed.
  \item Penalties are designed to \emph{increase} as \(d\) grows; late-detected misbehavior is more costly because it had more opportunity to pollute shared calibration signals.
  \item Validators who \emph{self-report} their misbehavior quickly incur a reduced penalty relative to those whose misbehavior is discovered by others at a later time.
\end{itemize}

Intuitively, this creates a \emph{temporal bond} between stake and behavior: validators are economically exposed to the risk that misbehavior will eventually be discovered, and they pay more if they attempt to conceal it longer.

\subsection{Temporal Slashing Function (Qualitative)}

We conceptualize a temporal slashing function as:
\begin{itemize}[leftmargin=*]
  \item A mapping from a normalized detection delay (for example, \(d\) rescaled by a maximum allowed delay \(T_{\max}\)) to a fraction of stake to be penalized.
  \item Strictly increasing in the detection delay: the penalty fraction for an event detected after a small lag is less than for one detected after a large lag.
  \item Bounded between a minimum penalty (for immediate detection or prompt self-report) and a maximum penalty (for very late or egregious misbehavior).
\end{itemize}

While the main text in this version intentionally omits an explicit formula, a typical design family would:
\begin{itemize}[leftmargin=*]
  \item Allow a gentle slope for very short delays (to tolerate benign noise or quick corrections).
  \item Introduce a steeper slope for moderate to long delays, so that attempts to ``hide'' misbehavior become quickly uneconomical.
\end{itemize}

\subsection{Self-Reporting Discount}

The model also includes a self-reporting path:
\begin{itemize}[leftmargin=*]
  \item If a validator detects an error in its own past contributions and self-reports within an explicit grace window (for example, within a small fixed interval after committing the event), the system applies a \emph{discounted} slashing rule.
  \item The discounted rule yields a smaller penalty fraction relative to the base rule that would apply if the misbehavior were instead discovered externally after a comparable delay.
\end{itemize}

Qualitatively, the design target is:
\begin{itemize}[leftmargin=*]
  \item For honest-but-noisy validators, self-reporting should be strictly beneficial in expectation: they lose less stake by promptly revealing errors than by leaving them to be discovered by others.
  \item For adversarial validators, the combination of temporal growth in penalties and the absence of self-reporting makes repeated misbehavior increasingly costly as detection delays grow.
\end{itemize}

\subsection{Contrast with Standard Slashing and CBFT}

In standard proof-of-stake and classical BFT-inspired systems, slashing schemes are primarily tuned to:
\begin{itemize}[leftmargin=*]
  \item deter equivocation and double-signing in consensus,
  \item penalize validator downtime or protocol violations that threaten safety or liveness.
\end{itemize}

Detection time is often treated as a secondary operational matter rather than a primary variable in the slashing function.
By contrast, in our setting:
\begin{itemize}[leftmargin=*]
  \item We use detection delay as a core component of governance penalties, because the \emph{duration} for which bad calibration signals remain uncorrected directly affects downstream decisions.
  \item The slashing model does not target consensus faults per se, but rather the integrity of shared FP calibration and other governance signals.
\end{itemize}

This shift in focus---from consensus safety to calibration integrity---motivates a different design space, in which temporal bonding and self-report incentives take center stage.

\section{Three-Phase Cold-Start Governance}
\label{sec:cold-start}

\subsection{Cold-Start Problem in Governance Networks}

In a cross-organizational calibration network, new validators (e.g., newly onboarded customers or partners) initially have no historical track record.
If all validators start with equal weight, early decisions may be disproportionately influenced by participants whose behavior has not yet been observed over enough events.

Na{\"i}ve solutions, such as simply assigning all nodes the same initial reputation and full weight, fail to:
\begin{itemize}[leftmargin=*]
  \item protect the network from early strategic misbehavior by new entrants, and
  \item reflect differences between long-standing participants with consistent behavior and new nodes with little history.
\end{itemize}

\subsection{Three-Phase Policy}

We propose a three-phase cold-start policy that modulates the weight of each validator based on the number of calibration submissions \(n\) it has contributed.
The policy is parameterized by two thresholds and a mapping from historical behavior to reputation, but the core structure is:

\paragraph{Phase 1: Zero-Weight Probation.}
\begin{itemize}[leftmargin=*]
  \item For a small initial number of submissions (illustratively, \(n \le 20\)), a validator carries \emph{zero} weight in aggregate calibration decisions.
  \item The system records these events and evaluates correctness over time, but they do not directly influence shared thresholds or governance rules.
\end{itemize}

\paragraph{Phase 2: Graduated Weight.}
\begin{itemize}[leftmargin=*]
  \item For an intermediate range of submissions (illustratively, \(20 < n \le 50\)), a validator is assigned a \emph{graduated} weight that increases as more events are observed.
  \item This weight is scaled by reputation, such that:
  \begin{itemize}
    \item validators exhibiting high correctness and prompt self-reporting gain influence more quickly, and
    \item validators showing misbehavior or noisy participation gain influence slowly or plateau at low weight.
  \end{itemize}
\end{itemize}

\paragraph{Phase 3: Full Weight.}
\begin{itemize}[leftmargin=*]
  \item After a sufficiently large number of submissions (illustratively, \(n > 50\)), a validator can reach \emph{full} weight, proportional to its reputation.
  \item In this phase, the validator behaves similarly to long-standing participants in terms of influence, but its reputation remains sensitive to ongoing behavior and slashing.
\end{itemize}

The threshold values (20 and 50) are deliberately presented as \emph{illustrative}, not normative.
An implementation may choose different thresholds based on volume, desired sensitivity, and risk tolerance.

\subsection{Illustrative Weight Profile}

Conceptually, the weight \(w_v\) for a validator \(v\) can be visualized as:
\begin{itemize}[leftmargin=*]
  \item \(w_v = 0\) for \(n \le 20\),
  \item \(w_v\) smoothly increasing from \(0\) to \(R_v\) as \(n\) grows from 20 to 50,
  \item \(w_v = R_v\) for \(n > 50\).
\end{itemize}

Here \(R_v\) is a separate reputation score that encodes correctness and temporal behavior (e.g., via an exponential moving average).
The exact functional shape (linear, concave, convex, or more complex) is left open to implementation.

\subsection{Relation to Reputation-Based Consensus}

Reputation-based consensus schemes have been proposed in the blockchain literature, often to:
\begin{itemize}[leftmargin=*]
  \item select block proposers or leaders based on past performance, and
  \item adjust voting power dynamically as behavior is observed.
\end{itemize}

Our three-phase cold-start regime is conceptually related but applied to \emph{governance of calibration signals} rather than block production:
\begin{itemize}[leftmargin=*]
  \item Phase 1 ensures that new validators cannot immediately affect shared thresholds.
  \item Phase 2 gradually onboards validators into governance, with influence gated by both history and reputation.
  \item Phase 3 places proven validators on equal footing, modulo their ongoing reputation and stake.
\end{itemize}

This distinction---reputation-based governance in a calibration network, not core blockchain consensus---is central to the prior art we intend to anchor.

\section{Application to FP Calibration and Model Governance}

\subsection{FP Calibration Network}

Consider a Phase Mirror-style environment where multiple customers and partners participate in a shared FP calibration network.
Each validator contributes observations about false positives and related governance signals for various rules:
\begin{itemize}[leftmargin=*]
  \item For a given rule, validators submit events indicating whether a given case is a false positive or not.
  \item The network aggregates these signals to tune global or segment-specific thresholds that drive automated or semi-automated decisions.
\end{itemize}

The temporal bonding and cold-start regime apply as follows:
\begin{itemize}[leftmargin=*]
  \item \textbf{Temporal stake bonding} governs how severely a validator is penalized if it systematically misreports FP status or otherwise manipulates calibration data, with penalties increasing in detection delay.
  \item \textbf{Self-reporting discounts} encourage validators to quickly correct mis-labeled FP data, reducing the cost of honest errors and aligning incentives toward transparency.
  \item \textbf{Three-phase cold-start} controls how new validators gain influence over shared thresholds, ensuring that early contributions are observed and evaluated before they meaningfully shift governance rules.
\end{itemize}

\subsection{Governance of Cross-Organizational Rules}

In this setting, the shared FP calibration network is both:
\begin{itemize}[leftmargin=*]
  \item a data aggregation layer (collecting FP labels and related signals), and
  \item a governance layer (deciding how those signals influence rules that affect all participants).
\end{itemize}

Temporal stake bonding and cold-start governance provide a mechanism for:
\begin{itemize}[leftmargin=*]
  \item limiting the influence of misbehaving organizations on calibration outcomes,
  \item rewarding validators that maintain high-quality, promptly corrected signals, and
  \item establishing a principled path for new participants to gain influence while being scrutinized.
\end{itemize}

This is distinct from payment or transaction ledgers; the core asset under governance is calibration integrity rather than monetary balances.

\section{Trade-Offs and Open Questions}

\subsection{Trust Recursion and Arbitration}

Any reputation-based system introduces a recursion: reputations govern decisions, and decisions in turn affect reputations.
This raises questions about:
\begin{itemize}[leftmargin=*]
  \item how disputes over misbehavior or correctness are resolved,
  \item how appeals are handled when validators contest slashing or reputation updates, and
  \item how to prevent collusion among validators to distort calibration signals or mislabel competitors.
\end{itemize}

Conceptually, one could:
\begin{itemize}[leftmargin=*]
  \item introduce an arbitration layer, inspired by systems like Kleros, where complex disputes are resolved by a separate pool of jurors or experts,
  \item require additional stake bonding for participation in arbitration, with its own temporal and cold-start profiles,
  \item and treat arbitration outcomes as inputs to reputation and slashing, rather than prescribing a fixed schedule here.
\end{itemize}

However, this paper deliberately refrains from fixing arbitration mechanics, juror selection rules, or related meta-governance structures.
Those are left as open design choices for specific deployments.

\subsection{Parameterization and Risk Profiles}

The design space for concrete parameter choices includes:
\begin{itemize}[leftmargin=*]
  \item the maximum detection delay \(T_{\max}\) (if explicitly specified),
  \item the shape of the temporal slashing function (e.g., linear vs. convex),
  \item the length of the self-report grace period and the magnitude of the self-report discount,
  \item the cold-start thresholds for submissions and the graduated weight schedule.
\end{itemize}

These choices depend on:
\begin{itemize}[leftmargin=*]
  \item typical event volume and latency in the calibration network,
  \item tolerance for noise versus adversarial manipulation,
  \item and legal or contractual constraints around penalties and dispute resolution.
\end{itemize}

We intentionally avoid committing to any specific numerical schedule beyond illustrative examples, so that deployments can tune these parameters without conflicting with the prior art or prematurely freezing a suboptimal profile.

\section{Deliberate Omissions and Boundary of this Disclosure}

To balance defensive publication with preservation of trade-secret space, this paper makes several deliberate omissions and non-claims.

\subsection{Stake Levels and Currency Units}

We consciously avoid:
\begin{itemize}[leftmargin=*]
  \item specifying exact stake levels or currency units (e.g., ``1{,}000 units of stake'' or denominating stake in any concrete asset),
  \item binding implementations to any particular mapping between subscription tiers and stake size.
\end{itemize}

Instead, we refer generically to stake \(S_v\) and to minimum stake requirements \(S_0\) (if any) as abstract quantities to be instantiated by specific deployments.

\subsection{Revenue-Sharing and Economic Schedules}

We also avoid:
\begin{itemize}[leftmargin=*]
  \item fixing subscription pricing or billing models for network participants,
  \item binding ourselves to any particular revenue-sharing fraction (e.g., ``30\% of slashing proceeds are sent to a treasury'' or similar),
  \item specifying how stake penalties are redistributed (burned, reallocated, or credited) beyond the abstract statement that they represent an economic cost.
\end{itemize}

These details are left to product design, financial modeling, and legal review for particular systems.

\subsection{Cryptographic and Protocol Details}

Finally, we do \emph{not} specify:
\begin{itemize}[leftmargin=*]
  \item identity, authentication, or key management schemes,
  \item transport protocols, message formats, or database layouts,
  \item cryptographic constructions such as zero-knowledge proofs, multi-party computation, or threshold signatures that may be layered on top of this governance model,
  \item deterministic encoding rules for numeric data, treatment of NaNs, NFC normalization, or duplicate-key handling.
\end{itemize}

Those aspects belong to implementation profiles and security specifications and can be defined separately, with their own versioning and anchoring.

\section{Conclusion}

We have presented a qualitative model of temporal stake bonding and a three-phase cold-start regime for governance of cross-organizational calibration signals.
The core prior art this paper intends to freeze is:
\begin{itemize}[leftmargin=*]
  \item the use of \emph{time-of-detection} as a first-class variable in stake slashing for governance of FP calibration and related signals,
  \item a three-phase (probation, graduated, full) cold-start policy for reputation-based aggregation of calibration signals in cross-organizational networks, and
  \item an explicit application of these ideas to a Phase Mirror-style FP calibration oracle, distinct from payment ledgers or core blockchain consensus.
\end{itemize}

By articulating this kernel and making the boundary of the disclosure explicit, we aim to establish defensible prior art for this class of mechanisms, while leaving substantial room for proprietary implementations, parameterization, and security engineering.

\appendix
\section{Mathematical Appendix}
\label{app:math}

This appendix collects formal statements and proofs of the key results used in the main text: the temporal lower bound on adversarial attack cost, the self-report incentive lemma, and an operator-norm view of the EWMA reputation update and its concentration properties. The focus is on a minimal, implementation-agnostic mathematical kernel.

\subsection{Preliminaries and Notation}

We restate the basic objects needed for the proofs.

\begin{itemize}[leftmargin=*]
  \item Validators \(V = \{v_1,\dots,v_n\}\) with stakes \(S_v \ge 0\).
  \item Discrete time \(t = 0,1,2,\dots\).
  \item Reputations \(R_v(t) \in [0,1]\) for each \(v\) and \(t\).
  \item Calibration events \(e = (v, r, \ell, t_0)\) with true labels \(\ell^*(e)\).
  \item For misbehavior events, detection time \(t_{\mathrm{detect}}(e)\) and delay
  \[
  d(e) = t_{\mathrm{detect}}(e) - t_0(e).
  \]
\end{itemize}

We assume:

\begin{assumption}[Bounded Detection Delay]
\label{ass:bounded-delay-app}
There exists \(T_{\max} \in (0,\infty)\) such that for all misbehavior events \(e\),
\[
0 \le d(e) \le T_{\max}.
\]
\end{assumption}

\begin{assumption}[Honest and Adversarial Types]
\label{ass:types-app}
There exist \(0 < \epsilon < \delta < 1\) such that for each validator \(v\):
\begin{itemize}[leftmargin=*]
  \item If \(v\) is honest, then \(\mathbb{P}[\ell = \ell^*(e)] \ge 1 - \epsilon\).
  \item If \(v\) is adversarial, then \(\mathbb{P}[\ell = \ell^*(e)] \le 1 - \delta\).
\end{itemize}
\end{assumption}

\begin{assumption}[Honest Majority]
\label{ass:honest-majority-app}
There exists \(\alpha > 1/2\) such that at least a fraction \(\alpha\) of validators (or of total stake) are honest.
\end{assumption}

We also recall the induced operator norm for a linear map \(T:\mathbb{R}^n \to \mathbb{R}^n\): with respect to a norm \(\|\cdot\|\),
\[
\|T\|_{\mathrm{op}} = \sup_{\bm{x} \neq 0} \frac{\|T\bm{x}\|}{\|\bm{x}\|}.
\]
In what follows, we use the Euclidean norm \(\|\cdot\|_2\) when referring to \(\|\cdot\|_{\mathrm{op}}\).

\subsection{Temporal Slashing and Attack Cost}

\subsubsection{Temporal Slashing Function}

\begin{definition}[Temporal Slashing Function]
\label{def:temporal-slash-app}
Let \(T_{\max} > 0\) satisfy Assumption~\ref{ass:bounded-delay-app}. Let
\[
f : [0,1] \to [0,1]
\]
be continuous and strictly increasing with
\[
f(0) = \beta_{\min}, \quad f(1) = \beta_{\max},
\]
where \(0 < \beta_{\min} \le \beta_{\max} \le 1\).
For a misbehavior event \(e\) by validator \(v\) with detection delay \(d(e) \in [0,T_{\max}]\), define the slashing penalty
\[
\mathrm{Slash}_v(e) = S_v \cdot f\!\left(\frac{d(e)}{T_{\max}}\right).
\]
\end{definition}

\begin{definition}[Self-Reporting Variant]
\label{def:self-report-app}
Let \(\tau \in (0,T_{\max}]\) and \(\gamma \in (0,1)\). For an event \(e\) self-reported within delay \(\tau\) by validator \(v\), define
\[
f_{\mathrm{sr}}(x) = \gamma f(x), \qquad x = \frac{d(e)}{T_{\max}},
\]
and a discounted slashing penalty
\[
\mathrm{Slash}^{\mathrm{sr}}_v(e) = S_v \cdot f_{\mathrm{sr}}\!\left(\frac{d(e)}{T_{\max}}\right).
\]
\end{definition}

\subsubsection{Lemma: Lower Bound on Attack Cost}

\begin{lemma}[Temporal Lower Bound on Attack Cost]
\label{lem:attack-cost-app}
Let \(v\) be an adversarial validator with stake \(S_v\). Suppose \(v\) performs \(k \ge 1\) misbehavior events \(e_1,\dots,e_k\), each detected within time \(T_{\max}\), so \(d(e_i)\in(0,T_{\max}]\). Assume:
\begin{enumerate}[leftmargin=*]
  \item There exists \(d_{\min} > 0\) such that \(d(e_i) \ge d_{\min}\) for all \(i\) (the adversary avoids immediate detection and does not self-report).
  \item Slashing uses the base function \(f\) from Definition~\ref{def:temporal-slash-app} (no self-report discount).
\end{enumerate}
Then the total slashing penalty satisfies
\[
\mathrm{Slash}^{\mathrm{tot}}_v
:= \sum_{i=1}^k \mathrm{Slash}_v(e_i)
= S_v \sum_{i=1}^k f\!\left(\frac{d(e_i)}{T_{\max}}\right)
\ge S_v \cdot k \cdot f\!\left(\frac{d_{\min}}{T_{\max}}\right).
\]
\end{lemma}

\begin{proof}
Since \(f\) is strictly increasing on \([0,1]\) and \(d(e_i) \ge d_{\min} > 0\) for each \(i\), we have
\[
f\!\left(\frac{d(e_i)}{T_{\max}}\right)
\ge f\!\left(\frac{d_{\min}}{T_{\max}}\right)
\]
for all \(i = 1,\dots,k\). Multiplying by \(S_v \ge 0\) yields
\[
\mathrm{Slash}_v(e_i)
= S_v f\!\left(\frac{d(e_i)}{T_{\max}}\right)
\ge S_v f\!\left(\frac{d_{\min}}{T_{\max}}\right).
\]
Summing over \(i\) gives
\[
\mathrm{Slash}^{\mathrm{tot}}_v
= \sum_{i=1}^k \mathrm{Slash}_v(e_i)
\ge \sum_{i=1}^k S_v f\!\left(\frac{d_{\min}}{T_{\max}}\right)
= S_v \cdot k \cdot f\!\left(\frac{d_{\min}}{T_{\max}}\right),
\]
which establishes the inequality.
\end{proof}

\paragraph{Interpretation.}
Even if the adversary ceases misbehavior after \(k\) events, its minimum cumulative loss scales linearly with \(k\), and the per-event loss is controlled by the smallest detection delay it can reliably achieve. The temporal dependence in \(f\) thus acts as a bond: attempts to delay detection increase the marginal cost of each additional attack.

\subsection{Self-Reporting Incentives}

\subsubsection{Mislabel and Self-Report Model}

\begin{definition}[Mislabel and Self-Report Model]
\label{def:mislabel-model-app}
Fix a validator \(v\) with mislabel probability \(\epsilon \in (0,1)\). For each event:
\begin{itemize}[leftmargin=*]
  \item With probability \(1-\epsilon\), the label is correct and no slashing occurs.
  \item With probability \(\epsilon\), the label is incorrect (misbehavior). Conditional on misbehavior:
  \begin{itemize}
    \item With probability \(p_{\mathrm{sr}} \in [0,1]\), the validator self-reports within delay \(\tau\).
    \item With probability \(1-p_{\mathrm{sr}}\), misbehavior is detected externally with normalized detection delay \(D \in (0,1]\).
  \end{itemize}
\end{itemize}
\end{definition}

\begin{definition}[Expected Loss per Misbehavior Event]
\label{def:loss-app}
Conditional on misbehavior, the expected loss to validator \(v\) under self-report probability \(p_{\mathrm{sr}}\) is
\[
L(p_{\mathrm{sr}})
= S_v\left[
  p_{\mathrm{sr}} f_{\mathrm{sr}}\!\left(\frac{\tau}{T_{\max}}\right)
  + (1-p_{\mathrm{sr}})\mathbb{E}[f(D)]
\right],
\]
where the expectation is over the distribution of \(D\).
\end{definition}

The unconditional expected loss per event is \(\epsilon L(p_{\mathrm{sr}})\); for incentive comparisons we can ignore the constant factor \(\epsilon\).

\subsubsection{Lemma: Self-Reporting Strictly Beneficial}

\begin{lemma}[Self-Reporting Strictly Beneficial]
\label{lem:self-report-app}
Assume:
\begin{enumerate}[leftmargin=*]
  \item \(f_{\mathrm{sr}}(x) = \gamma f(x)\) for some fixed \(\gamma \in (0,1)\).
  \item \(\tau/T_{\max} \in [0,1]\).
  \item The inequality
  \[
  \gamma f\!\left(\frac{\tau}{T_{\max}}\right) < \mathbb{E}[f(D)]
  \]
  holds for the distribution of normalized delays \(D\) for non-self-reported misbehavior.
\end{enumerate}
Then:
\begin{enumerate}[leftmargin=*]
  \item The expected loss \(L(p_{\mathrm{sr}})\) is strictly decreasing in \(p_{\mathrm{sr}}\in[0,1]\).
  \item Consequently, \(L(p_{\mathrm{sr}})\) is uniquely minimized at \(p_{\mathrm{sr}}^\star = 1\).
\end{enumerate}
\end{lemma}

\begin{proof}
Substituting \(f_{\mathrm{sr}}(x)=\gamma f(x)\) into Definition~\ref{def:loss-app} gives
\[
L(p_{\mathrm{sr}})
= S_v\left[
  p_{\mathrm{sr}} \gamma f\!\left(\frac{\tau}{T_{\max}}\right)
  + (1-p_{\mathrm{sr}})\mathbb{E}[f(D)]
\right].
\]
This expands to
\[
L(p_{\mathrm{sr}})
= S_v\left[
  \mathbb{E}[f(D)]
  + p_{\mathrm{sr}}\left(
    \gamma f\!\left(\frac{\tau}{T_{\max}}\right)
    - \mathbb{E}[f(D)]
  \right)
\right].
\]
Differentiating with respect to \(p_{\mathrm{sr}}\) yields
\[
\frac{dL}{dp_{\mathrm{sr}}}
= S_v\left[
  \gamma f\!\left(\frac{\tau}{T_{\max}}\right)
  - \mathbb{E}[f(D)]
\right].
\]
By assumption, the bracketed term is strictly negative, so \(\frac{dL}{dp_{\mathrm{sr}}} < 0\) for all \(p_{\mathrm{sr}}\in[0,1]\). Hence \(L\) is strictly decreasing and attains its unique minimum at the right endpoint \(p_{\mathrm{sr}} = 1\).
\end{proof}

\paragraph{Interpretation.}
By choosing \(\gamma\) and \(\tau\) so that the discounted penalty for quick self-report is below the expected penalty for non-self-reported misbehavior, the mechanism ensures that honest-but-noisy validators strictly prefer self-reporting in expectation. This aligns economic incentives with prompt error disclosure.

\subsection{EWMA Reputation as a Linear Operator}

We now analyze the EWMA reputation update as a linear operator on \(\mathbb{R}^n\), derive its operator norm, and sketch a concentration result for individual validators.

\subsubsection{Vector Formulation}

Stack reputations into a vector
\[
\bm{R}(t) = (R_1(t),\dots,R_n(t))^\top \in \mathbb{R}^n.
\]
Let \(\bm{\Delta}(t) = (\Delta_1(t),\dots,\Delta_n(t))^\top\), where \(\Delta_v(t) \in [0,1]\) is an outcome indicator (e.g., 1 for correct, 0 for misbehavior) at epoch \(t\).

The EWMA update is
\[
R_v(t+1) = \lambda R_v(t) + (1-\lambda)\Delta_v(t), \quad v = 1,\dots,n,
\]
with \(\lambda \in (0,1)\). In vector form,
\[
\bm{R}(t+1) = \lambda \bm{R}(t) + (1-\lambda)\bm{\Delta}(t).
\]

Define the linear map \(T:\mathbb{R}^n \to \mathbb{R}^n\) by \(T\bm{x} = \lambda \bm{x}\), and the affine update
\[
\bm{R}(t+1) = T\bm{R}(t) + (1-\lambda)\bm{\Delta}(t).
\]

\begin{lemma}[Operator Norm of EWMA Linear Part]
\label{lem:ewma-op-app}
Let \(T(\bm{x}) = \lambda \bm{x}\) with \(\lambda \in (0,1)\). With respect to the Euclidean norm \(\|\cdot\|_2\),
\[
\|T\|_{\mathrm{op}} = |\lambda|.
\]
In particular, \(\|T\|_{\mathrm{op}} = \lambda < 1\).
\end{lemma}

\begin{proof}
For any nonzero \(\bm{x} \in \mathbb{R}^n\),
\[
\frac{\|T\bm{x}\|_2}{\|\bm{x}\|_2}
= \frac{\|\lambda \bm{x}\|_2}{\|\bm{x}\|_2}
= \frac{|\lambda|\|\bm{x}\|_2}{\|\bm{x}\|_2}
= |\lambda|.
\]
Thus the ratio is constant in \(\bm{x}\), so taking the supremum over all nonzero \(\bm{x}\) yields \(\|T\|_{\mathrm{op}} = |\lambda|\).
\end{proof}

\paragraph{Consequence.}
Lemma~\ref{lem:ewma-op-app} shows that the homogeneous part of the EWMA recursion is a contraction when \(\lambda \in (0,1)\). Differences in initial reputations decay geometrically; the stochastic driving term \((1-\lambda)\bm{\Delta}(t)\) determines the long-run behavior.

\subsubsection{Closed-Form Representation}

\begin{lemma}[Unrolled EWMA Recursion]
\label{lem:ewma-unroll-app}
For any \(t \ge 1\),
\[
\bm{R}(t)
= \lambda^t \bm{R}(0)
  + (1-\lambda)\sum_{s=0}^{t-1} \lambda^{t-1-s}\bm{\Delta}(s).
\]
\end{lemma}

\begin{proof}
The result follows by induction. For \(t=1\),
\[
\bm{R}(1) = \lambda \bm{R}(0) + (1-\lambda)\bm{\Delta}(0),
\]
which matches the formula. Assume it holds for \(t\). Then:
\[
\begin{aligned}
\bm{R}(t+1)
&= \lambda \bm{R}(t) + (1-\lambda)\bm{\Delta}(t) \\
&= \lambda\left[
  \lambda^t \bm{R}(0)
  + (1-\lambda)\sum_{s=0}^{t-1}\lambda^{t-1-s}\bm{\Delta}(s)
\right]
+ (1-\lambda)\bm{\Delta}(t) \\
&= \lambda^{t+1}\bm{R}(0)
  + (1-\lambda)\sum_{s=0}^{t-1}\lambda^{t-s}\bm{\Delta}(s)
  + (1-\lambda)\bm{\Delta}(t) \\
&= \lambda^{t+1}\bm{R}(0)
  + (1-\lambda)\sum_{s=0}^{t}\lambda^{t-s}\bm{\Delta}(s),
\end{aligned}
\]
which is the claimed formula for \(t+1\).
\end{proof}

\subsubsection{Single-Validator Concentration}

For a single validator \(v\),
\[
R_v(t+1) = \lambda R_v(t) + (1-\lambda)\Delta_v(t),
\]
and by unrolling,
\[
R_v(t)
= \lambda^t R_v(0)
  + (1-\lambda)\sum_{s=0}^{t-1}\lambda^{t-1-s}\Delta_v(s).
\]

Assume \(\Delta_v(s)\in[0,1]\) and \(\mathbb{E}[\Delta_v(s)] = \mu_v\) for all \(s\), with \(\mu_v\in[0,1]\). Define
\[
m_v(t)
:= \lambda^t R_v(0)
  + (1-\lambda)\sum_{s=0}^{t-1}\lambda^{t-1-s}\mu_v,
\]
and the centered process
\[
X_v(t)
:= R_v(t) - m_v(t)
= (1-\lambda)\sum_{s=0}^{t-1}\lambda^{t-1-s}(\Delta_v(s) - \mu_v).
\]

\begin{lemma}[Concentration of EWMA Around Mean]
\label{lem:ewma-conc-app}
Fix a validator \(v\) and \(t \ge 1\). Assume:
\begin{enumerate}[leftmargin=*]
  \item \(\Delta_v(s) \in [0,1]\) for all \(s\).
  \item \(\Delta_v(s)\) are independent over \(s\) with \(\mathbb{E}[\Delta_v(s)] = \mu_v\).
\end{enumerate}
Then for any \(\eta > 0\),
\[
\mathbb{P}\big(|R_v(t) - m_v(t)| \ge \eta\big)
\le 2 \exp\!\left(
  -\frac{2\eta^2}{(1-\lambda)^2\sum_{s=0}^{t-1}\lambda^{2(t-1-s)}}
\right).
\]
\end{lemma}

\begin{proof}
Write
\[
R_v(t) - m_v(t)
= (1-\lambda)\sum_{s=0}^{t-1}\lambda^{t-1-s}(\Delta_v(s) - \mu_v)
= \sum_{s=0}^{t-1} a_s Y_s,
\]
where
\[
a_s = (1-\lambda)\lambda^{t-1-s}, \qquad
Y_s = \Delta_v(s) - \mu_v.
\]
Then \(Y_s \in [-\mu_v,1-\mu_v] \subseteq [-1,1]\) and \(\mathbb{E}[Y_s]=0\). By Hoeffding-type bounds for weighted sums of independent, bounded random variables, we have
\[
\mathbb{P}\left(\left|\sum_{s=0}^{t-1} a_s Y_s\right| \ge \eta\right)
\le 2\exp\left(
  -\frac{2\eta^2}{\sum_{s=0}^{t-1}(b_s - a_s')^2}
\right),
\]
where \([a_s',b_s]\) is a range containing \(a_sY_s\). Since \(|Y_s|\le 1\), we may take \(a_s'=-|a_s|\), \(b_s=|a_s|\), so \(b_s-a_s'=2|a_s|\) and
\[
\sum_{s=0}^{t-1}(b_s - a_s')^2
= \sum_{s=0}^{t-1}(2|a_s|)^2
= 4\sum_{s=0}^{t-1} a_s^2.
\]
Therefore
\[
\mathbb{P}(|R_v(t) - m_v(t)| \ge \eta)
\le 2 \exp\left(
  -\frac{2\eta^2}{4\sum_{s=0}^{t-1} a_s^2}
\right)
= 2\exp\left(
  -\frac{\eta^2}{2\sum_{s=0}^{t-1} a_s^2}
\right).
\]
Substituting \(a_s = (1-\lambda)\lambda^{t-1-s}\) gives
\[
\sum_{s=0}^{t-1} a_s^2
= (1-\lambda)^2\sum_{s=0}^{t-1}\lambda^{2(t-1-s)}
= (1-\lambda)^2\sum_{j=0}^{t-1}\lambda^{2j},
\]
which yields the bound in the statement (up to the standard constant factor; we present it in a form that exposes the geometric weighting explicitly).
\end{proof}

\paragraph{Gap Between Honest and Adversarial Reputations.}
Under Assumption~\ref{ass:types-app}, honest validators have \(\mu_v \ge 1-\epsilon\), while adversarial validators have \(\mu_v \le 1-\delta\) with \(\delta > \epsilon\). Lemma~\ref{lem:ewma-conc-app} shows that \(R_v(t)\) concentrates around \(m_v(t)\), which itself converges (for fixed \(\mu_v\)) to a value determined by \(\mu_v\). Thus, for sufficiently large \(t\), honest and adversarial validators occupy distinct intervals of typical reputation values with high probability, creating the separation used in the main text.

\subsection{Aggregation Operator and Phase-Dependent Weights}

Fix a time \(t\) and a rule \(r\). Let \(\bm{w}(t) = (w_1(t),\dots,w_n(t))^\top\) be the vector of phase-dependent weights, combining:
\begin{itemize}[leftmargin=*]
  \item a cold-start phase policy (probation, graduated, full) driven by the event count \(N_v(t)\), and
  \item reputation \(R_v(t)\) as described above.
\end{itemize}

Let \(\bm{x}(t) = (x_1(t),\dots,x_n(t))^\top\) be scalar reports (e.g., FP indicators or thresholds) from validators at time \(t\). Consider a linear aggregation operator
\[
A_t(\bm{x}(t)) = \sum_{v=1}^n w_v(t)x_v(t) = \bm{w}(t)^\top \bm{x}(t).
\]

\begin{lemma}[Operator Norm of Aggregation Map]
\label{lem:aggregation-op-app}
View \(A_t\) as a linear map from \((\mathbb{R}^n,\|\cdot\|_2)\) to \((\mathbb{R},|\cdot|)\). Then
\[
\|A_t\|_{\mathrm{op}} = \|\bm{w}(t)\|_2.
\]
\end{lemma}

\begin{proof}
For any nonzero \(\bm{x} \in \mathbb{R}^n\),
\[
\frac{|A_t(\bm{x})|}{\|\bm{x}\|_2}
= \frac{|\bm{w}(t)^\top \bm{x}|}{\|\bm{x}\|_2}.
\]
By Cauchy–Schwarz,
\[
|\bm{w}(t)^\top \bm{x}|
\le \|\bm{w}(t)\|_2\|\bm{x}\|_2,
\]
so the ratio is at most \(\|\bm{w}(t)\|_2\). Equality is attained when \(\bm{x}\) is parallel to \(\bm{w}(t)\), so the supremum is exactly \(\|\bm{w}(t)\|_2\).
\end{proof}

\paragraph{Interpretation.}
As time progresses and the cold-start regime moves validators from probation to full weight, the vector \(\bm{w}(t)\) becomes dominated by honest validators under Assumptions~\ref{ass:types-app} and~\ref{ass:honest-majority-app}. In operator-norm terms, the aggregation map \(A_t\) increasingly aligns with honest directions in \(\mathbb{R}^n\): its norm restricted to adversarial subspaces shrinks relative to that on honest subspaces, capturing in linear-algebraic language the asymptotic honest dominance described in the main text.

\bigskip

The results in this appendix isolate the mathematical content behind:
\begin{itemize}[leftmargin=*]
  \item temporal stake bonding and its linear lower bound on attack cost,
  \item incentive compatibility of self-reporting under discounted temporal penalties,
  \item contraction properties of the EWMA update operator and concentration of reputations around type-dependent means, and
  \item the operator-norm characterization of aggregate weighted decisions in the presence of phase-dependent weights.
\end{itemize}
They are deliberately stated in a way that is independent of particular parameter choices, detection mechanisms, or protocol implementations, so they can serve as a reusable analytical kernel for a class of governance and calibration systems.

\nocite{*}
\bibliographystyle{unsrt}  
\bibliography{references}  


\end{document}
